% ---------------------------------------------------------------------------
% Seção de resultados da próxima versão do artigo.
%
% Compila sozinho para revisão. Para entrar no artigo, mover o conteúdo entre
% \begin{document} e \end{document} e conferir que o preâmbulo de lá tem
% booktabs.
%
% Os números vêm de scripts/analisar_corrigido.py e scripts/contrastes.py,
% sobre as 1.723 execuções da campanha `conclusao`. Nenhum valor foi digitado
% à mão: as tabelas são as saídas dos scripts.
% ---------------------------------------------------------------------------

\documentclass[12pt,a4paper]{article}
\usepackage[utf8]{inputenc}
\usepackage[T1]{fontenc}
\usepackage[brazil]{babel}
\usepackage{booktabs}
\usepackage{geometry}
\geometry{margin=2.5cm}

\title{Quando a correção entra na conta do boato:\\
um artefato de medição em simulações de desinformação com agentes LLM}
\date{}

\begin{document}
\maketitle

\section{O defeito de medição}

O simulador classifica cada mensagem por uma linhagem operacional. A ação
\texttt{REFUTAR} inverte a marca dessa linhagem, e essa inversão é deliberada:
é o que permite que a correção seja entregue a quem já recebeu o boato, em vez
de ser barrada pela deduplicação junto com ele.

O problema é que a mesma marca era usada para classificar. Como ela alterna a
cada refutação, contestar uma correção devolve o texto à contagem do boato.
Nada no conteúdo retorna à alegação original: um texto que nega a alegação pela
segunda vez continua negando. Só a marca volta.

A correção separa as duas funções. A marca segue alternando, porque a entrega
depende disso, e a classificação passa a usar a profundidade de refutação, que
só cresce. A profundidade $0$ é o boato, a $1$ é a correção dele, e de $2$ em
diante a postura do texto não se resolve sem ler o conteúdo. Essa terceira faixa
recebe contagem própria, chamada aqui de linhagem ambígua.

\section{Reconstrução sobre os resultados gravados}

A auditoria de mensagens registra, por execução, cada publicação e cada
recebimento aceito, com o texto, o texto de origem e a marca de linhagem do
pai. Isso permite reconstruir a profundidade sem reexecutar a campanha.

A reconstrução foi aplicada às 1.723 execuções, que produziram 1.641 contrastes
pareados. Nenhuma execução ficou sem auditoria.

A validação é interna e não depende de confiança no procedimento: a
reconstrução também recalcula a contagem antiga, a partir dos mesmos eventos, e
compara com o valor gravado em cada execução. As duas coincidem em todas as
1.723. Divergência aqui indicaria leitura incorreta da auditoria, e portanto
invalidaria o alcance reconstruído.

\section{O tamanho do artefato}

A Tabela~\ref{tab:alcances} apresenta os alcances nos 82 controles. Os
intervalos vêm de reamostragem com 10.000 repetições, agrupada pelas 23
sementes estruturais, porque execuções que compartilham a semente compartilham
o grafo e a população inicial e não são realizações independentes.

\begin{table}[htbp]
\centering
\caption{Alcance nos controles, em porcentagem da população.}
\label{tab:alcances}
\begin{tabular}{lrrr}
\toprule
Alcance & Média & Mediana & IC 95\% \\
\midrule
Pela contagem antiga      & 36,20 & 2,0   & [26,60; 45,92] \\
Estrito do boato          &  6,80 & 2,0   & [2,49; 14,53] \\
Da correção               & 60,98 & 100,0 & [49,41; 73,68] \\
Ambíguo                   & 31,59 & 0,0   & [22,61; 41,15] \\
Limite superior do boato  & 36,20 & 2,0   & [26,60; 45,92] \\
\bottomrule
\end{tabular}
\end{table}

A diferença pareada entre as duas contagens, na mesma execução, tem média de
27,13 pontos percentuais, com intervalo de 21,47 a 33,00 sobre as 1.723
execuções. Em 85\% delas o boato não passou do emissor.

O alcance estrito é um limite inferior e o ambíguo não pode ser somado sem
leitura do conteúdo, então o alcance verdadeiro do boato está entre 6,80\% e
36,20\%. Fechar esse intervalo exige anotação de postura por revisores
independentes, que não faz parte destes resultados.

\section{Alcance por estímulo}

A Tabela~\ref{tab:estimulos} separa o resultado por texto de partida, e é onde
o artefato deixa de ser uma questão de magnitude.

\begin{table}[htbp]
\centering
\caption{Alcance por estímulo. A coluna ``só o emissor'' é a proporção de
execuções em que o boato não saiu do primeiro nó.}
\label{tab:estimulos}
\begin{tabular}{lrrrr}
\toprule
Estímulo & $n$ & Estrito (\%) & Antigo (\%) & Só o emissor \\
\midrule
\texttt{benefit\_auto\_a} & 395 & 5,57 & 13,21 & 0,84 \\
\texttt{benefit\_auto\_b} & 503 & 8,16 & 21,27 & 0,63 \\
\texttt{health\_dna\_a}   & 475 & 2,00 & 63,19 & 1,00 \\
\texttt{health\_dna\_b}   & 350 & 2,00 & 25,05 & 1,00 \\
\bottomrule
\end{tabular}
\end{table}

Nos dois estímulos de saúde, o alcance estrito é de 2,00\% com intervalo de
2,00 a 2,00, e a proporção de execuções em que o boato ficou no emissor é de
1,00 com intervalo de 1,00 a 1,00. Com 50 agentes, 2\% é exatamente o emissor.
Em nenhuma das 825 execuções desses dois estímulos a alegação falsa foi
retransmitida.

A contagem antiga, para o mesmo conjunto, reportava alcance médio de 63,19\% em
\texttt{health\_dna\_a}, com mediana de 94\%. Os 61 pontos percentuais de
diferença eram correções e contestações de correções.

Não se trata de um intervalo estreito, e sim de um intervalo degenerado: a
variância é nula porque o fenômeno medido não ocorreu. Os estímulos de economia
circularam de fato, e é deles que vêm os 4.303 textos distintos publicados na
linhagem do boato.

\section{Contrastes pareados}

A Tabela~\ref{tab:contrastes} compara cada tratamento ao controle do seu
próprio bloco, nas duas métricas lado a lado. Os intervalos são de 99,77\%,
correspondendo à correção de Bonferroni sobre as 22 condições declaradas.

Nas intervenções, o efeito é calculado apenas sobre as execuções em que o
tratamento foi de fato aplicado. Em $t=0$ isso vale para as 311 execuções. Em
$t=2$, 126 das 314 não receberam o tratamento porque a cascata se extinguiu
antes do segundo tick, 103 delas com uma única rodada. As linhas de $t=2$
descrevem portanto um subconjunto condicionado à sobrevivência da cascata, e
não devem ser lidas como efeito de aplicar o tratamento mais tarde.

\begin{table}[htbp]
\centering
\caption{Contrastes pareados por condição, nas duas métricas. Intervalos de
99,77\%, com correção de Bonferroni sobre as 22 condições.}
\label{tab:contrastes}
\footnotesize
\setlength{\tabcolsep}{4pt}
% Gerado por scripts/contrastes.py. Nao editar a mao.
% Intervalos de 99.77%, bootstrap agrupado pela seed estrutural.
\begin{tabular}{llrrrr}
\toprule
Família & Condição & $n$ & Seeds & Antigo (p.p.) & Estrito (p.p.) \\
\midrule
emitter & \texttt{cetico} & 78 & 22 & -6.95 [-26.44, 12.13] & -5.05 [-18.87, -0.14] \\
 & \texttt{fact\_checker} & 78 & 23 & 27.69 [5.35, 50.10] & -5.05 [-18.65, -0.18] \\
 & \texttt{indiferente} & 80 & 23 & -8.28 [-28.24, 12.54] & -4.92 [-18.09, -0.21] \\
 & \texttt{influenciador} & 77 & 22 & -8.99 [-27.06, 10.31] & -5.12 [-18.72, -0.14] \\
\midrule
intervention & \texttt{factcheckers\_hub\_t0} & 77 & 22 & 3.95 [-13.49, 20.77] & -4.00 [-16.49, 0.24] \\
 & \texttt{factcheckers\_hub\_t2} & 44 & 20 & 8.09 [-16.19, 36.77] & -4.00 [-19.36, 0.55] \\
 & \texttt{factcheckers\_peripheral\_t0} & 77 & 22 & 6.16 [-9.43, 24.06] & -1.53 [-6.16, 0.28] \\
 & \texttt{factcheckers\_peripheral\_t2} & 49 & 20 & 26.04 [-2.72, 54.51] & -2.00 [-10.45, 0.80] \\
 & \texttt{limit\_forwarding\_t0} & 78 & 22 & -25.82 [-41.03, -10.33] & -4.13 [-16.51, 0.09] \\
 & \texttt{limit\_forwarding\_t2} & 48 & 19 & -19.79 [-35.84, -0.61] & -3.12 [-13.89, 0.00] \\
 & \texttt{remove\_hubs\_t0} & 79 & 23 & -20.43 [-33.29, -5.97] & -4.13 [-16.75, -0.07] \\
 & \texttt{remove\_hubs\_t2} & 47 & 19 & -23.11 [-40.32, -5.07] & -5.23 [-26.68, 0.65] \\
\midrule
mutation & \texttt{literal\_forwarding} & 78 & 22 & -26.90 [-44.55, -9.54] & 0.23 [-1.15, 2.24] \\
\midrule
persistence & \texttt{double\_horizon} & 78 & 22 & 6.77 [-6.15, 21.32] & -0.49 [-1.32, 0.28] \\
\midrule
skeptics & \texttt{fraction\_0} & 78 & 22 & 0.03 [-18.38, 22.56] & 0.44 [-1.30, 3.73] \\
 & \texttt{fraction\_0.4} & 78 & 22 & 4.59 [-10.90, 23.04] & 0.05 [-0.48, 0.68] \\
 & \texttt{fraction\_0.6} & 79 & 22 & -6.13 [-26.24, 9.67] & -4.20 [-16.11, 0.00] \\
 & \texttt{fraction\_0.8} & 77 & 22 & -16.60 [-34.95, 0.58] & -4.05 [-16.75, 0.55] \\
 & \texttt{fraction\_1} & 78 & 22 & -12.33 [-33.79, 6.31] & -3.74 [-16.95, 0.39] \\
\midrule
topology & \texttt{barabasi\_albert} & 56 & 15 & -0.04 [-14.90, 13.00] & 1.46 [-0.95, 10.00] \\
 & \texttt{erdos\_renyi} & 57 & 15 & 4.60 [-14.45, 26.38] & -1.86 [-15.45, 7.84] \\
 & \texttt{watts\_strogatz} & 44 & 14 & 4.55 [-13.31, 23.72] & 4.27 [-1.18, 17.82] \\
\bottomrule
\end{tabular}

\end{table}

\section{O que muda nas conclusões}

A correção não altera apenas a magnitude dos efeitos. Em três pontos ela altera
a conclusão.

\paragraph{O efeito da mutação desaparece.} Proibir a alteração do texto
(\texttt{literal\_forwarding}) reduzia o alcance em 26,90 pontos percentuais
pela contagem antiga, com intervalo de $-44{,}55$ a $-9{,}54$, sobrevivendo à
correção de multiplicidade. Pela contagem estrita o efeito é de $+0{,}23$, com
intervalo de $-1{,}15$ a $2{,}24$. O que a proibição de alterar reduzia era a
circulação de correções reescritas, não a do boato.

\paragraph{Um efeito troca de sinal.} Com o emissor convertido em verificador,
a contagem antiga registra aumento de 27,69 pontos percentuais, com intervalo
de $5{,}35$ a $50{,}10$. A contagem estrita registra redução de 5,05, com
intervalo de $-18{,}65$ a $-0{,}18$. Os dois intervalos excluem o zero, em
direções opostas, e ambos sobrevivem à correção de multiplicidade. Um emissor
verificador produz correções, e eram elas que entravam na conta do boato.

Na métrica corrigida, os quatro perfis de emissor testados produzem efeitos
praticamente iguais, entre $-4{,}92$ e $-5{,}12$, todos com intervalo excluindo
o zero. Nessas condições o que importa não é qual perfil substitui o emissor
padrão, e sim deixar de usar o amplificador. A diferenciação entre perfis que a
contagem antiga sugeria vinha de quanto texto corretivo cada um produzia.

\paragraph{O limite de retransmissão deixa de se sustentar.} Pela contagem
antiga, limitar retransmissões em $t=0$ reduzia o alcance em 25,82 pontos, com
intervalo de $-41{,}03$ a $-10{,}33$. Pela contagem estrita o efeito é de 4,13,
com intervalo de $-16{,}51$ a $0{,}09$, que inclui o zero. A única intervenção
de rede cujo efeito sobrevive na métrica corrigida é a remoção de hubs em
$t=0$, com 4,13 pontos e intervalo de $-16{,}75$ a $-0{,}07$, e por margem
estreita.

O padrão aparente de limiar na fração de céticos também não sobrevive. Pela
contagem antiga os efeitos crescem em módulo de 0,03 a $-16{,}60$ conforme a
fração sobe de $0$ a $0{,}8$. Pela contagem estrita nenhuma das cinco frações
tem intervalo que exclua o zero.

\section{Limitações}

Os resultados descrevem as condições executadas e não estimam comportamento de
pessoas. A população é sintética e não foi calibrada; os quatro estímulos são
fixos, de modo que não há base para generalizar por tema.

A separação entre alcance estrito e ambíguo é operacional. Ela indica onde a
leitura do conteúdo é necessária, e não substitui essa leitura. Enquanto a
anotação de postura não for feita, o alcance do boato permanece como intervalo,
e não como ponto.

A contagem de agentes convencidos usa a mesma marca de linhagem para decidir a
direção da atualização de crença, e portanto carrega o mesmo defeito. Esse caso
não pode ser corrigido sobre os resultados gravados, porque a crença evolui
durante a execução, e exigiria reexecutar a campanha. A métrica não é usada
como desfecho neste trabalho.

As linhas de $t=2$ dependem da sobrevivência da cascata, como já observado. A
comparação direta entre momentos de aplicação exigiria um desenho que trate
essa seleção, o que não foi feito.

Por fim, a ablação de cache e de regra de difusão não foi executada, e a
comparação quantitativa com modelos SIR e SEIR continua pendente. A decomposição
das fontes de variação entre rede, perfis, ativação e modelo exige desenho
fatorial próprio: o simulador registra sementes separadas para cada uma, mas a
campanha não as variou de forma independente.

\end{document}
